\documentclass[11pt,a4paper]{article}
\usepackage[utf8]{inputenc}
\usepackage[T1]{fontenc}
\usepackage[brazil]{babel}
\usepackage{lmodern}
\usepackage{microtype}
\usepackage{geometry}
\geometry{left=2.3cm,right=2.3cm,top=2.2cm,bottom=2.2cm}
\usepackage{booktabs}
\usepackage{tabularx}
\usepackage{array}
\usepackage{graphicx}
\usepackage{float}
\usepackage{xcolor}
\usepackage{hyperref}
\usepackage{enumitem}
\usepackage{tikz}
\usetikzlibrary{arrows.meta,positioning}
\usepackage{fancyhdr}
\usepackage{titlesec}
\usepackage{setspace}

\definecolor{azul}{HTML}{173B57}
\definecolor{cinza}{HTML}{F1F3F5}
\hypersetup{colorlinks=true,linkcolor=azul,urlcolor=azul}
\titleformat{\section}{\Large\bfseries\color{azul}}{}{0pt}{}
\titleformat{\subsection}{\large\bfseries\color{azul}}{}{0pt}{}
\setlength{\parindent}{1.1cm}
\setlength{\parskip}{0.25em}
\setlist[itemize]{topsep=0.25em,itemsep=0.15em}
\setlist[enumerate]{topsep=0.25em,itemsep=0.15em}
\renewcommand{\arraystretch}{1.16}

\pagestyle{fancy}
\fancyhf{}
\lhead{Simulador didático de sistemas operacionais}
\rhead{Sistemas Operacionais}
\cfoot{\thepage}
\setlength{\headheight}{14pt}

\begin{document}

\begin{titlepage}
\thispagestyle{empty}
\centering
\vspace*{1.5cm}
{\Large Disciplina de Sistemas Operacionais\par}
\vspace{1.2cm}
{\Huge\bfseries\color{azul} Simulador Didático de Sistemas Operacionais\par}
\vspace{0.4cm}
{\LARGE Relatório técnico\par}
\vspace{1.4cm}

\begin{tabular}{ll}
Anthony Kewin Pinto & RA 115738\\
Isadora Vischi da Cruz & RA 114640\\
Lucas Gabriel Risso & RA 114544\\
Victor Tazoi & RA 114923\\
Yan Vinícius Pimenta Silva & RA 115273\\
\end{tabular}

\vfill
\textbf{Professor:} Maurício Acconcia Dias\\[0.4cm]
\textbf{Repositório:} \url{https://github.com/grupo-fho/simulador-so}\\[0.4cm]
\textbf{Versão documentada:} Java 21\\[0.4cm]
20 de setembro de 2026
\end{titlepage}

\section{Objetivo do projeto}

O trabalho consiste em um simulador de mecanismos de sistemas operacionais. A proposta não é substituir o sistema operacional da máquina, mas representar de forma controlada situações estudadas na disciplina: criação e execução de processos e threads, escalonamento da CPU, paginação, operações de entrada e saída e acesso a um sistema de arquivos simplificado.

O tempo é representado por um clock lógico. Dessa forma, o resultado não depende da velocidade do computador em que o programa é executado. Uma mesma carga, com os mesmos parâmetros, produz a mesma sequência de eventos e permite comparar políticas sob as mesmas condições.

A versão atual oferece uma CLI, uma interface gráfica em Swing e geração de logs e arquivos CSV. A interface gráfica é uma camada adicional; o núcleo utilizado por ela é o mesmo da execução por linha de comando.

\section{Arquitetura}

A implementação foi dividida por responsabilidade. O núcleo da simulação fica em \texttt{Simulador}, enquanto as regras de CPU, memória, E/S e arquivos ficam em classes próprias. Essa separação facilita a leitura e evita que toda a lógica fique concentrada em uma única classe.

\begin{figure}[H]
\centering
\begin{tikzpicture}[
box/.style={draw=azul, rounded corners, fill=cinza, minimum width=3.3cm, minimum height=0.9cm, align=center},
arrow/.style={-{Latex[length=2mm]}, thick, draw=azul},
node distance=1.2cm and 1.2cm
]
\node[box] (main) {Main / Configuração / Carga};
\node[box, below=of main] (sim) {Simulador + Evento};
\node[box, below left=of sim, xshift=-1.5cm] (esc) {Escalonador};
\node[box, below=of sim, xshift=-1.7cm] (mem) {Memória};
\node[box, below=of sim, xshift=1.7cm] (io) {Dispositivos};
\node[box, below right=of sim, xshift=1.5cm] (arq) {Arquivos};
\draw[arrow] (main) -- (sim);
\draw[arrow] (sim) -- (esc);
\draw[arrow] (sim) -- (mem);
\draw[arrow] (sim) -- (io);
\draw[arrow] (sim) -- (arq);
\end{tikzpicture}
\caption{Organização geral do simulador.}
\end{figure}

\texttt{Processo} representa as informações equivalentes ao PCB. A classe mantém identificador, prioridade, chegada, estado agregado, threads, espaço de endereçamento e descritores de arquivos. \texttt{ThreadSimulada} representa o TCB e armazena estado, contador de programa, registradores simulados, pilha lógica e dados usados nas métricas.

O diagrama de classes completo está em \texttt{docs/diagrama-classes.mmd}.

\section{Clock, eventos e estados}

O simulador avança o clock até o próximo evento agendado. Quando mais de um evento ocorre no mesmo instante, a ordem adotada é: conclusão de E/S ou paginação, chegada, término de unidade da CPU e fim de troca de contexto. Empates dentro do mesmo tipo seguem a sequência de inserção.

As threads percorrem os estados NOVO, PRONTO, EXECUTANDO, BLOQUEADO e FINALIZADO. Cada mudança é registrada no log com a causa que provocou a transição. Por exemplo, uma falta de página leva a thread de EXECUTANDO para BLOQUEADO; quando a página é carregada, ela volta para PRONTO.

O estado do processo é obtido a partir das suas threads e das operações pendentes. Por isso, o PCB não precisa repetir exatamente todas as transições do TCB.

\section{Escalonamento da CPU}

Foram implementadas as três políticas obrigatórias: FCFS, Round Robin e Prioridade. No FCFS, a thread permanece na CPU até bloquear ou terminar. No Round Robin, o quantum limita o tempo contínuo de execução; se o surto não terminar, a thread retorna ao fim da fila. Na política de Prioridade, o menor valor numérico representa maior prioridade.

A política de Prioridade é não preemptiva. Portanto, a chegada de uma thread com prioridade melhor não interrompe imediatamente a que já está executando. A nova thread será considerada no próximo despacho.

A troca de contexto possui custo configurável. Esse intervalo entra no clock total, mas é separado do tempo de CPU útil. O simulador também conta quantas trocas foram realizadas.

\section{Gerenciamento de memória}

A memória usa paginação com páginas e molduras de tamanho fixo. Cada processo possui sua própria tabela de páginas. Um endereço lógico é dividido em número da página e deslocamento; quando a página está presente, o endereço físico é calculado pela moldura associada.

Se a página não estiver residente, ocorre uma falta de página. A thread é bloqueada e o pedido entra na fila do paginador. Havendo moldura livre, ela é utilizada. Caso contrário, a vítima é escolhida por FIFO. Depois do tempo de atendimento, a tabela de páginas é atualizada e a thread volta para a fila de prontas.

A referência que gerou a falta é concluída no evento de paginação. Isso evita registrar a mesma referência duas vezes quando a thread retoma a execução.

\section{Entrada e saída}

O simulador possui um dispositivo de bloco e um de caractere, cada um com fila FCFS e tempo de serviço próprio. Uma operação bloqueante retira a thread da CPU até a conclusão. A conclusão gera uma interrupção simulada e devolve a thread ao estado PRONTO.

Também existe operação não bloqueante. Nesse caso, a thread continua executando e o resultado é entregue posteriormente pela caixa de entrada do TCB. Não é utilizada espera ativa.

\section{Sistema de arquivos}

O sistema de arquivos é hierárquico e inicia em um diretório raiz. Arquivos e diretórios possuem metadados simplificados, como nome, tipo, permissões, tamanho e instantes lógicos de criação e modificação.

Foram implementadas as operações de criar, abrir, ler, escrever, fechar, remover e listar. Cada processo possui uma tabela local de descritores, ligada à tabela global de arquivos abertos.

Os dados dos arquivos usam alocação encadeada em blocos. Essa estratégia não exige blocos contíguos, mas o acesso a posições mais distantes depende do percurso da cadeia. O volume é mantido apenas durante a execução e não substitui o sistema de arquivos real do computador.

\section{Interface e execução}

A carga de trabalho fica fora do código-fonte. O arquivo descreve processos, chegadas, prioridades, threads e instruções. Parâmetros de escalonamento, memória e dispositivos podem ser informados pela CLI ou pela interface gráfica.

A GUI foi feita em Swing. Ela utiliza \texttt{Main.executar}, o mesmo fluxo da linha de comando, e roda a simulação em \texttt{SwingWorker} para manter a janela responsiva. O uso de \texttt{SwingWorker} não cria paralelismo entre as threads simuladas; o núcleo continua orientado a eventos e controlado pelo clock lógico.

\section{Métricas}

As métricas são calculadas a partir do clock lógico. As definições utilizadas são:

\begin{itemize}
\item \textbf{Tempo de retorno}: finalização menos chegada;
\item \textbf{Tempo de espera}: soma dos intervalos na fila de prontas;
\item \textbf{Tempo de resposta}: chegada até a primeira utilização da CPU;
\item \textbf{Utilização da CPU}: unidades úteis divididas pelo tempo total;
\item \textbf{Utilização dos dispositivos}: tempo ocupado dividido pelo tempo total;
\item \textbf{Throughput}: processos concluídos por unidade de tempo;
\item \textbf{Taxa de faltas}: faltas de página divididas pelo número de referências;
\item \textbf{Trocas de contexto}: quantidade de trocas e sobrecarga acumulada.
\end{itemize}

\section{Experimentos}

Os experimentos foram executados com cargas controladas. Nas comparações, os parâmetros que não estavam sendo estudados foram mantidos constantes.

\subsection{Comparação das políticas de CPU}

A primeira comparação usa a mesma carga mista para FCFS, Round Robin com quantum 3 e Prioridade.

\begin{table}[H]
\centering
\small
\begin{tabular}{lrrrrr}
\toprule
Política & Total & Trocas & Resposta média & Espera média & Retorno médio\\
\midrule
FCFS & 70 & 13 & 17,75 & 30,25 & 63,33\\
RR, q=3 & 87 & 24 & 6,25 & 42,00 & 69,00\\
Prioridade & 75 & 12 & 15,00 & 28,75 & 50,33\\
\bottomrule
\end{tabular}
\caption{Comparação das políticas de escalonamento.}
\end{table}

Nesse cenário, o Round Robin oferece a menor resposta média, pois distribui a CPU mais cedo entre as threads. Em contrapartida, realiza mais trocas de contexto e aumenta o tempo total. A política de Prioridade apresenta o menor retorno médio entre os três resultados, enquanto o FCFS termina o conjunto em menos tempo que o RR. As métricas mostram que cada política favorece um aspecto diferente da execução.

A quantidade de faltas também muda entre as políticas, mesmo com a configuração de memória igual. Isso ocorre porque o escalonamento altera a ordem em que as referências são feitas.

\subsection{Efeito do quantum}

O segundo experimento usa Round Robin e altera apenas o quantum.

\begin{table}[H]
\centering
\small
\begin{tabular}{rrrrr}
\toprule
Quantum & Total & Trocas & Resposta média & CPU útil\\
\midrule
1 & 122 & 54 & 3,25 & 44,26\%\\
3 & 87 & 24 & 6,25 & 62,07\%\\
6 & 76 & 15 & 10,00 & 71,05\%\\
\bottomrule
\end{tabular}
\caption{Efeito do quantum no Round Robin.}
\end{table}

O quantum 1 reduz o tempo de resposta, mas causa 54 trocas para as mesmas 54 unidades de CPU útil. O quantum 6 reduz a sobrecarga e o tempo total, mas demora mais para oferecer a primeira oportunidade de CPU. O resultado mostra o compromisso entre responsividade e custo de troca.

\subsection{Quantidade de molduras}

A terceira comparação usa uma carga de memória que repete três páginas e altera somente a quantidade de molduras.

\begin{table}[H]
\centering
\small
\begin{tabular}{rrrrr}
\toprule
Molduras & Faltas & Taxa de faltas & Total & Ocioso\\
\midrule
2 & 9 & 100,00\% & 65 & 45\\
3 & 3 & 33,33\% & 29 & 15\\
\bottomrule
\end{tabular}
\caption{Comparação da quantidade de molduras.}
\end{table}

Com três molduras, as três primeiras referências carregam as páginas e as seis seguintes encontram seus dados residentes. Com duas molduras, a sequência provoca substituições sucessivas. Nessa carga, aumentar uma moldura reduz de forma clara os bloqueios e o tempo total.

\section{Testes e validação}

A versão atual foi validada com Java 21. A suíte registra 191 verificações aprovadas e cobre regras do núcleo, cenários de integração, interface e validações do lançador do DOOM. Entre os pontos verificados estão transições de estado, filas de escalonamento, quantum, métricas, tradução de endereços, FIFO, E/S bloqueante e não bloqueante, descritores e operações de arquivo.

Além dos testes automatizados, o repositório possui cargas pequenas que permitem conferir manualmente a ordem dos eventos. Essa conferência é útil para demonstrar que o resultado do simulador pode ser acompanhado passo a passo.

Antes da entrega, a distribuição deve ser testada novamente em uma cópia limpa, seguindo apenas as instruções do README.

\section{Limitações}

O simulador trabalha com uma única CPU lógica e tempos de serviço fixos. Não foram implementados multiprocessamento real, TLB, páginas sujas, envelhecimento de prioridade, deadlocks, semáforos, cache de blocos, journaling ou persistência física do volume simulado.

O paginador possui seu próprio atendimento e não converte automaticamente uma falta de página em requisição ao objeto que representa o disco. Da mesma forma, as operações do sistema de arquivos não são transformadas em E/S do dispositivo de bloco. Essas escolhas simplificam o modelo e devem ser consideradas ao interpretar os resultados.

A integração com Chocolate Doom também é separada do núcleo. O jogo é iniciado por \texttt{ProcessBuilder} como processo real do sistema hospedeiro e não é escalonado pelo simulador.

\section{Conclusão}

A implementação integra os mecanismos principais pedidos no trabalho em um mesmo eixo de tempo lógico. O uso de logs e métricas permite acompanhar as decisões do simulador e comparar configurações sem depender do tempo real da máquina.

Os experimentos mostram três efeitos importantes: a política de CPU altera resposta, espera e sobrecarga; o quantum do Round Robin modifica diretamente o número de trocas; e a quantidade de molduras influencia as faltas de página e o tempo de execução. Esses resultados reforçam o objetivo didático do projeto, que é observar como decisões de um sistema operacional afetam o comportamento da carga.

\section{Ferramentas e materiais}

O desenvolvimento e a validação utilizam JDK 21, scripts PowerShell, Python para automação dos experimentos e Git/GitHub para versionamento. O enunciado da disciplina é a referência para o escopo e os critérios de entrega. Ferramentas auxiliares adicionais que tenham sido utilizadas pelo grupo devem permanecer declaradas no README conforme as regras acadêmicas da disciplina.

\end{document}
